\documentclass[../Main.tex]{subfiles}

\externaldocument{../Main}

\begin{document}
\chapter{Foundations of Probability and Measure}

Probability theory provides a mathematical framework for studying random phenomena. Many processes in nature, science, engineering, and finance involve outcomes that cannot be predicted with certainty. Examples include the outcome of a coin toss, the arrival time of customers in a queue, or the evolution of asset prices in financial markets. Probability theory allows us to model such situations and to reason quantitatively about uncertainty.

The modern formulation of probability theory is based on \textbf{measure theory}. In measure theory, the goal is to assign a non-negative number to certain subsets of a given set. This number is called a \textbf{measure} and represents a generalized notion of size, extending familiar ideas such as length, area, or volume. A fundamental property of a measure is \textbf{countable additivity}: if a collection of sets is disjoint, the measure of their union is the sum of their measures.

For mathematical reasons, it is not possible in general to assign a measure to every subset of a set. Instead, one restricts attention to a special collection of subsets called a \sfield. A set equipped with a \sfield is called a \textbf{measurable space}, and measures are defined only on these measurable sets.

Probability theory can be viewed as a special case of measure theory, where we are assigning a measure that represents our belief of occurrence of events. This measure is called a \textbf{Probability measure} and it has the additional property that the total measure of the entire space is equal to one. The triplet $(\probspace)$ is called a \textbf{probability space} where : 

\begin{itemize}
    \item The set $\Omega$ is the \textbf{sample space}. It contains all possible outcomes of the experiment. Once an outcome $\omega \in \Omega$ is fixed, all random quantities associated with the experiment are completely determined.
    \item The collection $\A$ is a \sfield of subsets of $\Omega$, called \textbf{events}. These are the sets to which probabilities can be assigned. From the point of view of measure theory, they are precisely the measurable sets. An event $A \in \A$ represents the occurrence of a given property of the outcome $\omega$.
    \item The function $\PP : \A \to [0,1]$ is a \textbf{probability measure}. For each event $A \in \A$, the value $\PP(A)$ represents the probability that the event $A$ occurs.
\end{itemize}

This axiomatic formulation of probability theory was introduced by \textbf{Andrey Kolmogorov in 1933}. The Kolmogorov axioms provide a simple and powerful foundation for probability and allow many results to be developed in a unified and rigorous way.

Although probability theory relies on measure theory, the goals of the two subjects are different. Measure theory studies general notions of size and integration, while probability theory focuses on modeling and analyzing randomness. In this chapter, we mainly develop probability theory while referring to measure-theoretic concepts whenever they are needed.

A prior knowledge of measure theory can be helpful but is not strictly required. The reader is assumed to have seen basic notions such as \sfields, measurable functions, and measures. Whenever these ideas are used, we briefly recall their meaning and explain their role in probability.

The aim of this chapter is to present the fundamental concepts of probability theory in a clear and structured way, combining mathematical rigor with intuitive interpretations based on random experiments.

\medskip
\textbf{Important remark.}
In many situations, the explicit form of the probability measure $\PP$ will not be specified. Instead, we focus on the properties of measurable functions defined on $(\measspace)$, which are called \emph{random variables}.

\medskip
\textbf{Terminology.}
As in measure theory, sets of measure zero play an important role in probability theory.  A property depending on $\omega \in \Omega$ is said to hold \emph{almost surely} (a.s.) if it holds for all $\omega$ in an event of probability one. Thus, \textit{almost surely} has the same meaning as "almost everywhere" in measure theory.

\medskip
\textbf{Notation.}
For a random variable $\bX$, the notation $\PP(\bX \in A)$ means :
\[ \PP(\{\omega \in \Omega : \bX(\omega) \in A\}). \]
\newpage


\section{General Definitions}
In this section, we introduce the fundamental objects of probability theory. We begin with the definition of a probability space, which provides the mathematical framework for modeling random experiments. We then define random variables, which allow us to describe and analyze numerical quantities generated by randomness.

These definitions form the basis of all subsequent results in this chapter. They are stated in a general and rigorous setting, with close connections to measure theory whenever necessary.



\subsection{Probability space}
\defnr{Probability space }{probability_space}{
A \textbf{probability space} is a triple $(\probspace)$, where:
\begin{itemize}
    \item $\Omega$ is a non-empty set called the \textbf{sample space},
    \item $\A$ is a \sfield of subsets of $\Omega$, and
    \item $\PP: \A \to [0,1]$ is a \textbf{probability measure}.
\end{itemize}
}
We remind the reader that a \sfield $\A$ is a collection of subsets of $\Omega$ such that:
\begin{itemize}
    \item $\Omega \in \A$,
    \item if $A \in \A$ then $A^c = \Omega \setminus A \in \A$,
    \item if $A_1, A_2, \ldots \in \A$ then $\cup_{i=1}^\infty A_i \in \A$.
\end{itemize}

\defnr{Probability measure }{probability_measure}{
A \textbf{probability measure} $\PP$ on a \sfield $\A$ of a set $\Omega$ is a function $\PP: \A \to [0,1]$ satisfying:
\begin{itemize}
    \item $\PP(\Omega) = 1$
    \item For any countable sequence of disjoint sets $A_1, A_2, \ldots \in \A$,
    \[
    \PP\left( \bigcup_{i=1}^\infty A_i \right) = \sum_{i=1}^\infty \PP(A_i)
    \]
\end{itemize}
}

\proppr{measure_properties}{
We have the following results :
    \begin{enumerate}
        \item If $A \subset B$ then : \[ \PP(B\setminus A) = \PP(B) - \PP(A) \]
        \item If $A,B \in \A$  then : \[ \PP(A \cup B) = \PP(A) + \PP(B) - \PP(A \cap B) \]
        \item If $A_n \in \A$ and $A_n \subset A_{n+1}$ for every $n\in\NN$ then : \[ \PP(\bigcup_{n\in\NN} A_n ) = \lim_{\goto{n}} \PP(A_n)\]
        \item If $A_n \in \A$ and $A_{n+1} \subset A_{n}$ for every $n\in\NN$ then : \[ \PP(\bigcap_{n\in\NN} A_n ) = \lim_{\goto{n}} \PP(A_n)\]
        \item If $A_n \in \A$  for every $n\in\NN$ then : \[ \PP(\bigcup_{n\in\NN} A_n ) \le \sum_{n\in\NN} \PP(A_n)\]
    \end{enumerate}
}{ 
    Left exercise for the reader.
}

\subsection{Monotone class}
The axioms defining a $\sigma$-field are often stronger than those required in many arguments. In practice, one frequently begins with a simpler collection of sets, such as intervals or rectangles, and seeks to extend properties established on this collection to the $\sigma$-field it generates. This motivates the study of weaker systems whose closure properties are sufficient for such extension arguments. Among the most important are monotone classes (also called $\lambda$-systems or Dynkin systems). In this section, we focus on monotone classes, which are obtained by replacing the requirement of closure under arbitrary countable unions in the definition of a $\sigma$-field (Definition~\ref{defn:probability_space}) with the weaker requirement of closure under countable disjoint unions. We conclude by proving Dynkin's Monotone Class Theorem (also known as the $\pi$--$\lambda$ theorem), a fundamental result that allows properties established on a $\pi$-system to be extended to the $\sigma$-field it generates.

\defnr{Monotone class}{Monotone_class}{
Let $\Omega$ be a set and let $\M \subset \P(\Omega)$.   The collection $\M$ is called a \emph{monotone class} if:
\begin{enumerate}
    \item $\Omega \in \M$;
    \item if $A,B \in \M$ and $A \subset B$, then $B \setminus A \in \M$;
    \item if $(A_n)_{n \in \NN}$ is an increasing sequence in $\M$, i.e.
    \[ A_1 \subset A_2 \subset \cdots, \]
    then
    \[ \bigcup_{n \in \NN} A_n \in \M.\]
\end{enumerate}
}
Every $\sigma$-algebra is a monotone class. The converse is false in general. However, a monotone class is a $\sigma$-algebra if and only if it is closed under finite intersections. Indeed, closure under finite intersections yields closure under complements and finite unions, while the monotone class property then extends finite unions to countable unions, giving the defining properties of a $\sigma$-algebra.

It is therefore useful to introduce the notion of a $\pi$-system, namely a collection of sets closed under finite intersections. With this terminology, a $\sigma$-algebra may be viewed as the combination of a $\pi$-system and a $\lambda$-system (equivalently, a monotone class), that is,
\[
\boxed{\sigma\text{-algebra}=\pi\text{-system}+\lambda\text{-system}.}
\]
This observation underlies many uniqueness arguments in measure theory, most notably the $\pi$-$\lambda$ theorem.



Arbitrary intersections of monotone classes are again monotone classes. Given a collection $\C \subset \P(\Omega)$, we define the \emph{monotone class generated by $\C$}, denoted $\M(\C)$, as the intersection of all monotone classes containing $\C$.

\thmpr{Monotone Class Theorem}{monotone_class_theorem}{
    Let $\C \subset \P(\Omega)$ be a collection of sets that is closed under finite intersections. Then
    \[\M(\C) = \sigma(\C).\]
    Consequently, if $\M$ is any monotone class such that $\C \subset \M$, then
    \[\sigma(\C) \subset \M.\]
}{
It is enough to prove the first assertion. Since any \sfield is a monotone class, it is clear that $\M(C) \subset \sigma(\C)$. To prove the reverse inclusion, it is enough to verify that $\M(\C)$ is a \sfield. As explained above, it suffices to verify that $\M(\C)$ is closed under finite intersections.

For every $A \in \P(\Omega)$, set
\[ \M_A = \{B \in \M(\C) : A \cap B \in \M(\C)\}. \]

Fix $A \in \C$. Since $\C$ is closed under finite intersections, we have $\C \subset \M_A$. Let us verify that $\M_A$ is a monotone class:

\begin{itemize}
\item $\Omega \in \M_A$ is trivial.
\item If $B, B' \in \M_A$ and $B \subset B'$, then $A \cap (B' \setminus B) = (A \cap B') \setminus (A \cap B) \in \M(\C)$ and thus $B' \setminus B \in \M_A$.
\item If $B_n \in \M_A$ for every $n \geq 0$ and the sequence $(B_n)_{n \geq 0}$ is increasing, we have $A \cap (\cup_{n \in \NN} B_n) = \cup_{n \in \NN} (A \cap B_n) \in \M(\C)$ and therefore $\cup_{n \in \NN} B_n \in \M_A$.
\end{itemize}

Since $\M_A$ is a monotone class that contains $\C$, $\M_A$ also contains $\M(\C)$. We have thus obtained that
\[
\forall A \in \C, \ \forall B \in \M(\C), \ A \cap B \in \M(\C).
\]

This is not yet the desired result, but we can use the same idea one more time. Precisely, we now fix $A \in \M(\C)$. According to the first part of the proof, $\C \subset \M_A$. By exactly the same arguments as in the first part of the proof, we get that $\M_A$ is a monotone class. It follows that $\M(\C) \subset \M_A$, which shows that $\M(\C)$ is closed under finite intersections, and completes the proof.
}

Although the Monotone Class Theorem is easy to state, its true importance lies in the proof strategy that it provides. In practice, one rarely works directly with the monotone class $\M(\C)$. Instead, one defines the collection of all sets for which a desired property holds, proves that this collection forms a monotone class containing a suitable generating algebra, and then invokes the Monotone Class Theorem to extend the property to the entire generated $\sigma$-field. This four-step argument appears repeatedly throughout probability theory and measure theory. For future reference, the general scheme is summarized in Figure~\ref{fig:monotone_class_strategy}.

\begin{figure}[ht]
\centering
\begin{tikzpicture}[
    node distance=1.0cm,
    >=Latex,
    every node/.style={
        draw,
        rounded corners=2pt,
        align=center,
        text width=15cm,
        minimum height=1cm,
        font=\small
    },
    every edge/.style={draw, thick, ->}
]

\node (step1) {
\textbf{Step 1. Choose a generating algebra $\C$.}\\[2mm]
Verify that the desired property $P(A)$ holds for every $A\in\C$.
};

\node (step2) [below=of step1] {
\textbf{Step 2. Define the good class}
\[
\G=\{A\in\sigma(\C): P(A)\text{ holds}\}.
\]
Show that $\C\subseteq\G$.
};

\node (step3) [below=of step2] {
\textbf{Step 3. Show that $\G$ is a monotone class.}\\[2mm]
Verify that the property is preserved under increasing unions and
decreasing intersections.
};

\node (step4) [below=of step3] {
\textbf{Step 4. Apply the Monotone Class Theorem}
\[
\sigma(\C)\subseteq\G.
\]
Hence, the property $P(A)$ holds for every measurable set
$A\in\sigma(\C)$.
};

\draw (step1) -- (step2);
\draw (step2) -- (step3);
\draw (step3) -- (step4);

\end{tikzpicture}
\caption{The standard proof strategy for applying the Monotone Class Theorem.}
\label{fig:monotone_class_strategy}
\end{figure}

Here are some examples : 
\exm{Uniqueness of probability measures}{
Suppose $\PP$ and $\QQ$ are probability measures on $(\Omega,\A)$ satisfying
\[
\PP(A)=\QQ(A),\qquad A\in\C,
\]
where $\C$ is an algebra generating $\A$. To prove that $\PP=\QQ$, proceed as follows.

\begin{enumerate}
    \item Define the collection of \emph{good} sets:
    \[
    \G=\{A\in\A:\PP(A)=\QQ(A)\}.
    \]

    \item By assumption,
    \[
    \C\subseteq\G.
    \]

    \item Show that $\G$ is a monotone class. This follows from the continuity of finite measures from below and above.

    \item Since $\C$ is an algebra generating $\A$, the Monotone Class Theorem yields
    \[
    \A=\sigma(\C)\subseteq\G.
    \]
    Hence,
    \[
    \PP(A)=\QQ(A),\qquad A\in\A,
    \]
    so that $\PP=\QQ$.
\end{enumerate}
}

\exm{Extension of an integral identity}{
Suppose $f,g\ge0$ are measurable functions satisfying
\[
\int_A f\,d\mu=\int_A g\,d\mu,
\qquad A\in\C,
\]
where $\C$ is an algebra generating $\A$. To extend this identity to all measurable sets, proceed as follows.

\begin{enumerate}
    \item Define
    \[
    \G=\left\{A\in\A:\int_A f\,d\mu=\int_A g\,d\mu\right\}.
    \]

    \item Clearly,
    \[
    \C\subseteq\G.
    \]

    \item Show that $\G$ is a monotone class. This follows from the Monotone Convergence Theorem.

    \item The Monotone Class Theorem implies
    \[
    \A=\sigma(\C)\subseteq\G,
    \]
    so the identity holds for every measurable set.
\end{enumerate}
}

\exm{Equality of localized expectations}{
Let $X$ and $Y$ be integrable random variables such that
\[
\EE[X\1_A]=\EE[Y\1_A],
\qquad A\in\C,
\]
where $\C$ is an algebra generating $\A$. To prove that the equality holds for every measurable set, proceed as follows.

\begin{enumerate}
    \item Define
    \[
    \G=\left\{A\in\A:\EE[X\1_A]=\EE[Y\1_A]\right\}.
    \]

    \item By hypothesis,
    \[
    \C\subseteq\G.
    \]

    \item Verify that $\G$ is a monotone class. This follows from the Dominated Convergence Theorem.

    \item Applying the Monotone Class Theorem gives
    \[
    \A=\sigma(\C)\subseteq\G,
    \]
    and therefore
    \[
    \EE[X\1_A]=\EE[Y\1_A],\qquad A\in\A.
    \]
\end{enumerate}
}

\subsection{Random variables}
\defnr{Random variable}{random_variable}{
A function from $\Omega$ to $\RR^n$ is called a random variable (in measure theory is called measurable function), meaning an application $\bX$ such that:
\[ \bX^{-1}(B) \in \A, \text{ for each set } B \in \B(\RR^n) \]
where $\B(\RR^n)$ denotes the Borel \footnote{The Borel \sfield on $\RR^n$ is generated by the open subsets of $\RR^n$.  The Lebesgue measure on $\RR^n$ can be first defined on rectangles of the form $[a_1,b_1]\times\cdots\times[a_n,b_n]$ by $(b_1-a_1)\cdots(b_n-a_n)$, and then uniquely extended to a measure on the Borel \sfield. For more details, see the chapter \parencite[chap.~3]{Probability}.} \sfield. }

Based on this definition, we can define a probability measure in the probabilistic space $(\RR^n, \B(\RR^n), \PP_\bX)$, where the probability measure (also called the distribution of $\bX$) $\PP_\bX$ is defined as follows:
\[ \PP_\bX(B) = \PP(\bX^{-1}(B)) = \PP(\{ \bX \in B \}), \text{ for each set } B \in \B(\RR^n) \]
We call $\PP_\bX$ the law (or the distribution) of the random variable $\bX$. In order to characterize the distribution of a random variable, we use cumulative distribution functions and probability density functions.

\defnr{Cumulative distribution function}{cumulative_distribution_function}{
Let $\bX$ be a random variable with value space $\RR^n$. For each $\bx \in \RR^n$, we define the \textbf{cumulative distribution function} as follows:
\[ F_\bX(\bx) = \PP(\bX \leq \bx) =   \PP(X_1 \leq x_1, \cdots, X_n \leq x_n)  \]
}

\defnr{Probability density function}{density_function}{
Let $\bX$ be a random variable with value space $\RR^n$. For each $\bx \in \RR^n$, we define the \textbf{probability density function} as follows:
\[ f_\bX(\bx) = \frac{\partial^n}{\partial \bx} F_\bX(\bx) = \frac{\partial^n}{\partial x_1 \cdots \partial x_n} F_\bX(\bx) \]
if it exists on a non-negligible set of $\RR^n$. This function is a density for the probability measure $\PP_\bX$ such that:
\[ \PP_\bX(B) = \int_B d\PP_\bX =  \int_B f_\bX(\bx) d\lambda(\bx) \]
Where $\lambda$ is the Lebesgue measure on $\B(\RR^n)$
}
Finally, we have the following theorem, which ensures that a constructed variable is indeed a random variable.

\thmpr{Operations on random variables}{random_variables_operations}{
\begin{enumerate}
    \item Let $X_1,X_2$ be two random variables on $\RR$, and $f : \RR^2 \mapsto \RR$ a measurable function (w.r.t Lebesgue measure and Borel \sfield), then  $(X_1,X_2), \quad f(X_1,X_2)$ are random variables.
    \item Let $(X_n)$ be a sequence of random variables that map into $\bar{\RR}$ then :
    \[ \sup _{n \in \NN} X_n, \quad \inf_{n \in \NN} X_n, \quad \limsup _{\goto{n}} X_n, \quad  \liminf_{\goto{n}} X_n \]
    are also random variables. In particular, if the sequence $(X_n)$  converges pointwise, its limit is a random variable. 
    \item Let $(X_n)$ be a sequence of random variables that map into $\bar{\RR}$ then :
    The set $A$ of all $x \in \Omega$ such that $X_n(\omega)$ converges in $\RR$ as $\goto{n}$ is measurable ($A\in\A$). Furthermore, the function $h: \Omega \longrightarrow \RR$ defined by
    \[
    Y(\omega):= \begin{cases}
        \lim _{\goto{n}} X_n(\omega) & \text { if } \omega \in A, \\
        0 & \text { if } \omega \notin A,
                \end{cases}
    \]
    is a random variable.
\end{enumerate}
}{
The proofs could be found in \parencite[sec.~1.3]{Probability}. 
}
Finally we have this two lemmas, that would be used frequently in most of probability-propositions proofs.


\lempr{Approximation by Simple Random Variables}{approximation_simple_random_variables}{
Let \( X \) be a non-negative random variable with values in $\RR$. Then there is a sequence \( (\varphi_n )\) of simple random variables with \(\varphi_{n+1} \geq \varphi_n\) such that \( X = \lim \varphi_n \) at each point of \( \Omega \).
Where a simple random variable is a random variable having the form : 
\[
    X = \sum_{i=1}^m {c}_i \1_{E_i}
\]
For some integer $m$ and real positive values $c_i$. 
}{
The proof is left exercise for reader with the hint to use : 
\[
E_{n,k} = \{ x : k2^{-n} \leq X(\omega) < (k + 1)2^{-n} \}, \text{ and the sequence } \varphi_n = 2^{-n} \sum_{k=0}^{2^{2n}} k \1_{E_{n,k}}.
\]
}

\lempr{Doob-Dynkin Lemma}{Doob_Dynkin}{
If \( \bY\) with values in $\RR^m$ is \( \sigma(\bX) \)-measurable where $\bX$ is a random variable with values in $\RR^n$, there is a measurable function \( f : \RR^n \to \RR^m \) such that \( \bY= f( \bX) \).
}{
We'll proof the theorem for real valued random variables (with values in $\RR$). 
\begin{enumerate}

    \item \textbf{ Step 1: Indicator Functions} 
    First, consider the case where \( Y = \1_A \) for some \( A \in \sigma(X) \). 
    Since \( A \in \sigma(X) \), by definition, there exists a Borel set \( B \in \B(\RR) \) such that \( A = X^{-1}(B) \).
    Define the function \( f : \RR \to \RR \) by
    \[ f(x) = \1_B(x). \]
    Then, for any \( \omega \in \Omega \),
    \[ Y(\omega) = \1_A(\omega) = \1_{\bX^{-1}(B)}(\omega) = \1_B(X(\omega)) = f(X(\omega)). \]
    Hence, \( Y = f(X) \) for this case.

    \item \textbf{ Step 2: Simple Functions}
    Now, consider \( Y \) as a simple function, i.e., \( Y = \sum_{i=1}^n c_i \1_{A_i} \) where \( c_i \in \RR \) and \( A_i \in \sigma(X) \).
    For each \( A_i \in \sigma(X) \), there exists a Borel set \( B_i \in \B(\RR) \) such that \( A_i = X^{-1}(B_i) \).
    Define \( f : \RR \to \RR \) by
    \[ f(x) = \sum_{i=1}^n c_i \1_{B_i}(x). \]
    Then, for any \( \omega \in \Omega \),
    \[ Y(\omega) = \sum_{i=1}^n c_i \1_{A_i}(\omega) = \sum_{i=1}^n c_i \1_{\bX^{-1}(B_i)}(\omega) = \sum_{i=1}^n c_i \1_{B_i}(X(\omega)) = f(X(\omega)). \]
    Thus, \( Y = f(X) \) for this case.

    \item \textbf{ Step 3: Positive Measurable Functions}
    Consider \( Y \) as a positive \( \sigma(X) \)-measurable function. There exists a sequence of simple functions \( Y_n \) such that \( Y_n \nearrow Y \). (use of the lemma \ref{lem:approximation_simple_random_variables})
    From Step 2, for each \( n \), there exists a Borel-measurable function \( f_n : \RR \to \RR \) such that
    \[ Y_n = f_n(X). \]
    Since \( Y_n \nearrow Y \), we have \( f_n(X) \to Y \) pointwise almost everywhere. 
    Define \( f(x) = \lim_{\goto{n}} f_n(x) \). Since each \( f_n \) is Borel-measurable, and the pointwise limit of Borel-measurable functions is Borel-measurable, \( f \) is Borel-measurable.
    Therefore, \( Y = f(X) \).

    \item \textbf{ Step 4: General Measurable Functions}
    Finally, consider \( Y \) as an arbitrary \( \sigma(X) \)-measurable function. We can decompose \( Y \) into its positive and negative parts:
    \[ Y = Y^+ - Y^-, \]
    where \( Y^+ = \max(Y, 0) \) and \( Y^- = -\min(Y, 0) \).
    Both \( Y^+ \) and \( Y^- \) are positive \( \sigma(X) \)-measurable functions. By Step 3, there exist Borel-measurable functions \( f^+ \) and \( f^- \) such that
    \[ Y^+ = f^+(X) \quad \text{and} \quad Y^- = f^-(X). \]
    Define \( f : \RR \to \RR \) by
    \[ f(x) = f^+(x) - f^-(x). \]
    Then,
    \[ Y = Y^+ - Y^- = f^+(X) - f^-(X) = f(X). \]
    Thus, \( Y = f(X) \) for some Borel-measurable function \( f \).
\end{enumerate} }
A large number of proofs in probability theory follow the same general strategy. The idea is to first establish a result for very simple random variables and then extend it step by step to more general cases. This approach reflects the way many objects in probability theory are constructed from simpler ones.

More precisely, proofs often proceed through the following steps:

\begin{figure}[H]

\centering
\begin{tikzpicture}[
    node distance=1.1cm,
    >=Latex,
    every node/.style={
        draw,
        rounded corners=3pt,
        align=center,
        text width=15cm,
        minimum height=1.15cm,
        font=\small
    },
    every edge/.style={draw, thick, ->}
]

\node (s1) {
\textbf{Step 1: Indicator Functions}\\[2mm]
Prove the result for random variables of the form
\[
Y=\1_A,\qquad A\in\A.
\]
};

\node (s2) [below=of s1] {
\textbf{Step 2: Simple Random Variables}\\[2mm]
Extend the result by linearity to finite linear combinations
of indicator functions,
\[
Y=\sum_{i=1}^{n} a_i\1_{A_i}.
\]
};

\node (s3) [below=of s2] {
\textbf{Step 3: Positive Measurable Random Variables}\\[2mm]
Approximate by an increasing sequence of simple functions
\[ Y_n\uparrow Y, \]
and use the approximation lemma.
};

\node (s4) [below=of s3] {
\textbf{Step 4: General Measurable Random Variables}\\[2mm]
Write
\[ Y=Y^{+}-Y^{-}, \]
apply Step~3 to $Y^{+}$ and $Y^{-}$ separately,
and combine the results.
};

\node (final) [below=of s4] {
\textbf{Conclusion}\\[2mm]
The desired statement now holds for all measurable
random variables.
};

\draw (s1) -- (s2);
\draw (s2) -- (s3);
\draw (s3) -- (s4);
\draw (s4) -- (final);

\end{tikzpicture}
\caption{The standard four-step proof strategy used throughout probability theory.}
\label{fig:approximation_proof_strategy}
\end{figure}


In many situations, some of these steps may be omitted when they follow directly from previous arguments. Nevertheless, this four-step method provides a common structure that appears repeatedly in probability theory.

\section{Expectation and Moments of Random Variables}
\subsection{Expectation}
\defnr{Expectation}{Expectation}{
Let $\bX$ be a random variable with value space $\RR^n$. The \textbf{expectation} or \textbf{expected value} of a random variable $\bX$ with respect to a probability measure $\PP$ is defined as
\[
\EE[\bX] = \int_\Omega \bX \, d\PP =  \int_{\RR^n} \bx \, d\PP_\bX = \int_{\RR^n} \bx f_\bX(\bx) d\bx
\]
The integral here is a Lebesgue integral\footnote{The Lebesgue integral is the generalisation of integration on the measurable spaces, for the rigorous definition check the chapter \parencite[chap.~2]{Probability}} with respect to the probability measure. We require the integral to exist, and the last equality holds only if a density exists. }

We recall some properties of expectations as follow :
Let \( X \) and \( Y \) be random variables, and let \( a \) and \( b \) be constants. The main properties of expectations are:
\begin{enumerate}
    \item \textbf{Probability representation:}
    \begin{equation*}
    \PP[X\in A] = \EE(\1_{X\in A})
    \end{equation*}

    \item \textbf{Linearity:}
    \begin{equation*}
    \EE[a\bX+ b\bY] = a\EE[\bX] + b\EE[\bY]
    \end{equation*}

    \item \textbf{Expectation of a Constant:}
    \begin{equation*}
    \EE[a] = a
    \end{equation*}

    \item \textbf{Non-negativity:}
    If \( \bX\geq 0 \) almost surely, then
    \begin{equation*}
    \EE[\bX] \geq 0
    \end{equation*}

    \item \textbf{Finiteness:}
    \[ \EE[X] < \infty, \Longrightarrow  X < \infty \quad \PP\text{-a.s.} \]

    \item \textbf{Zero Expectation:}
    If \( X \) is a nonnegative random variable and
    \[ \EE[X] = 0 \;\Longleftrightarrow\; X = 0 \quad \PP\text{-a.s.} \]

    \item \textbf{Equality of Expectations:}
    If $Y$ is another nonnegative random variable and
    \[ X = Y \quad \PP\text{-a.s.} \]
    then
    \[ \EE[X] = \EE[Y]. \]

    \item \textbf{Law of the Unconscious Statistician:}
    If \( g \) is a measurable function and \( \bX\) is a random variable, then
    \begin{equation*}
    \EE[g(\bX)] = \int g(\bX)d\PP 
    \end{equation*}

    \item \textbf{Fatou's Lemma:}
    If \( (X_n) \) is a sequence of positive real random variables, then
    \begin{equation*}
    \EE[\liminf_{\goto{n}} X_n] \leq \liminf_{\goto{n}} \EE[X_n]
    \end{equation*}

    \item \textbf{Monotone Convergence Theorem:}
    If \( (X_n) \) is a sequence of positive real random variables that converges to \(X\) almost surely and the sequence is a.s. increasing, then
    \begin{equation*}
    \EE[\lim_{\goto{n}} X_n] = \lim_{\goto{n}} \EE[X_n]
    \end{equation*}

    \item \textbf{Dominated Convergence Theorem:}
    If \( (X_n) \) is a sequence of real random variables that converges to \( X\) almost surely and is dominated by an integrable random variable \( Y \) (i.e., \( |X_n| \leq Y \) almost surely for all \( n \)), then
    \begin{equation*}
    \EE[\lim_{\goto{n}} X_n] = \lim_{\goto{n}} \EE[X_n]
    \end{equation*}
\end{enumerate}

The last three theorems are fundamental results from measure theory and
Lebesgue integration. We do not present their proofs here, since most
applications in probability only require knowing when they can be applied.
They are the principal tools for exchanging limits and expectations.

A good rule of thumb is to proceed in the following order whenever you
encounter an expression of the form
\[ \EE\!\left[\lim_{n\to\infty}X_n\right]. \]

\begin{enumerate}
    \item Check whether \( X_n \) is increasing and \(X_n\rightarrow X\) almost surely\footnote{this is saying that limite of $X_n$ is $X$ on a set $A$ of $\omega\in A$ with $\PP(A)=1$}.
    If so, apply the \textbf{Monotone Convergence Theorem}.

    \item Otherwise, check whether there exists an integrable random variable
    \(Y\) satisfying
    \[
    |X_n|\leq Y
    \quad\text{a.s. for every }n.
    \]
    If so, apply the \textbf{Dominated Convergence Theorem}.

    \item If neither of the above hypotheses is satisfied, but the random
    variables are nonnegative, try using \textbf{Fatou's Lemma}. Although it
    only provides an inequality,
    \[
    \EE[\liminf X_n]
    \leq
    \liminf \EE[X_n],
    \]
    it is often sufficient to obtain bounds or prove convergence results.
\end{enumerate}

Another tool from real analysis is the Stone-Weierstrass theorem which is generally paired with the fact that :

$\EE[f(X)]=\EE[f(Y)]$ for every \textbf{bounded continuous} function $f: \RR\to \RR$. Then $\PP_X = \PP_Y$.

\thmr{Stone-Weierstrass theorem }{Stone_Weierstrass}{
Let \(K \subseteq \RR\) be a compact set (for example, \(K=[a,b]\)). Let 
\[ A \subseteq C(K,\CC)\]
be a subalgebra, where \(C(K,\CC)\) denotes the space of all continuous
complex-valued functions on \(K\), equipped with the uniform norm
\[ \|f\|_{\infty}=\sup_{x\in K}|f(x)|. \]

Assume that:

\begin{enumerate}
    \item \(A\) is a subalgebra of \(C(K,\CC)\); that is, whenever
    \(f,g\in A\) and \(\alpha,\beta\in\CC\),
    \[ \alpha f+\beta g\in A, \qquad fg\in A.     \]

    \item The constant function \(1\) belongs to \(A\).

    \item \(A\) separates points of \(K\); that is, for every distinct \(x,y\in K\), there exists \(f\in A\) such that
    \[ f(x)\neq f(y). \]
\end{enumerate}

Then \(A\) is dense in \(C(K,\CC)\) with respect to the uniform norm. Equivalently,
\[ \overline{A}^{\|\cdot\|_{\infty}} = C(K,\CC). \]

That is, for every \(f\in C(K,\CC)\) and every \(\varepsilon>0\), there exists
\(g\in A\) such that
\[\sup_{x\in K}|f(x)-g(x)|<\varepsilon. \]
}




Before completing the section we highly invit the reader to consult the section of $\L^p$ spaces \ref{sec:Lp spaces}


\subsection{Variance}
\defnr{Variance -- Covariance -- Correlation}{second_moments}{
Let $X \in L_\RR^2(\probspace)$ be a real-valued random variable.
\begin{enumerate}
\item \textbf{Variance.}  
The \textbf{variance} of the random variable $X$ measures how much $X$ varies around its mean. It is defined by
\[ \mathrm{Var}(X) = \EE\big[(X - \EE[X])^2\big]. \]
The \textbf{standard deviation} of $X$ is defined as
\[ \sigma(X) = \sqrt{\mathrm{Var}(X)}. \]

\item \textbf{Covariance.}  
Let $X$ and $Y$ be two random variables in $L^2(\probspace)$.  
The \textbf{covariance} of $X$ and $Y$ measures how the two variables vary together and is defined by
\[ \mathrm{Cov}(X,Y) = \EE\big[(X - \EE[X])(Y - \EE[Y])\big]
= \EE[XY] - \EE[X]\EE[Y]. \]

\item \textbf{Correlation.}  
The \textbf{correlation} between $X$ and $Y$ is a normalized version of the covariance and is defined by
\[ \rho(X,Y) = \frac{\mathrm{Cov}(X,Y)}{\sigma(X)\sigma(Y)}. \]
It takes values between $-1$ and $1$.

\item \textbf{Covariance matrix.}  
Let $\bX = (X_1,\dots,X_n)$ be a random vector in $\RR^n$ with $\bX \in L_{\RR^n}^2(\probspace)$. The \textbf{covariance matrix} of $\bX$ is defined by
\[
\mathrm{Var}(\bX) 
= \EE\big[(\bX - \EE[\bX])(\bX - \EE[\bX])^T\big]
= \big(\mathrm{Cov}(X_i,X_j)\big)_{1 \le i,j \le n}.
\]
This matrix belongs to $M_n(\RR)$.
\end{enumerate}
}

\thmpr{Markov Inequalities}{Markov_Inequalities}{
\textbf{Markov's Inequality} If $X$ is a nonnegative random variable and $a>0$,
\[ \PP(X \geq a) \leq \frac{1}{a} \EE[X] \] 
\textbf{Bienaymé-Chebyshev Inequality} If $X \in L^2(\probspace)$ and $a>0$,
\[\PP(|X-\EE[X]| \geq a) \leq \frac{1}{a^2} \operatorname{var}(X) \]
                }{
    A classical proof, left exercise for the reader.\\
    Hint : Use $\EE[\1_{X \geq a}] \leq \EE[X/a]$ for Markov's inequality, and apply Markov's inequality to the nonnegative random variable $(X-\EE[X])^2$ for Chebyshev's inequality.
                }

\subsection{Characteristic Function}
We start with a basic definition. We use the notation $x \cdot y$ for the standard Euclidean scalar product in $\RR^d$.

\defnr{Characteristic function}{Characteristic_function}{
Let $X$ be a random variable with values in $\RR^d$. The characteristic function of $X$ is the function $\Phi_X : \RR^d \longrightarrow \CC$ defined by
\[ \Phi_X(\xi) := \EE[\exp(i \xi \cdot X)], \qquad \xi \in \RR^d. \]
}

The definition of expectation \ref{defn:Expectation} allows us to write
\[ \Phi_X(\xi) = \int_{\RR^d} e^{i \xi \cdot x} \,d\PP_X(x) \]

and we can also view $\Phi_X$ as the Fourier transform of the law $\PP_X$ of $X$. For this reason, we sometimes write $\Phi_X(\xi) = \widehat{\PP_X}(\xi)$. It's easy to notice (in the case $d = 1$), the dominated convergence theorem ensures that the function $\Phi_X$ is (bounded and) continuous on $\RR^d$.

Our first goal is to show that the characteristic function determines the law $\PP_X$ of $X$ (as the name suggests!). This is equivalent to showing that the Fourier transform acting on probability measures on $\RR^d$ is injective. We begin with an important calculation for a special case of the Gaussian distribution. We will study this case in more detail in Section \ref{sec:gaussian_sec}, where we will derive the following formula:
\[ \Phi_X(\xi)=\exp \left(\mathrm{i} m \xi-\frac{\sigma^2 \xi^2}{2}\right), \quad \xi \in \RR \]

\thmpr{Characteristic function determines the law}{law_determined_by_characteristic_func}{
The characteristic function of a random variable $X$ with values in $\RR^d$ determines the law of this random variable. In other words, the Fourier transform acting on probability measures on $\RR^d$ is injective.
}{
Let us consider the case $d = 1$. For every $\sigma > 0$, let $g_\sigma$ be the density of the Gaussian $\N(0,\sigma^2)$ distribution:
\[
g_\sigma(x) = \frac{1}{\sigma \sqrt{2\pi}} \exp\left(-\frac{x^2}{2\sigma^2}\right), \qquad x \in \RR.
\]

Let $\mu$ be a probability measure on $\RR$, and set
\[
\begin{cases}
f_\sigma(x) &= \int_{\RR} g_\sigma(x - y)\,\mu(dy) \stackrel{(\mathrm{def})}{=} g_\sigma * \mu(x), \\ 
d\mu_\sigma(x) &= f_\sigma(x)\,dx.
\end{cases}
\]

Let $C_b(\RR)$ denote the space of all bounded continuous functions $\varphi : \RR \longrightarrow \RR$, and recall that a probability measure $\nu$ on $\RR$ is characterized by the values of $\int \varphi(x)\,d\nu(x)$ for $\varphi \in C_b(\RR)$. The result of the theorem (when $d = 1$) will follow if we are able to prove that:
\begin{enumerate}
\item $\mu_\sigma$ is determined by $\widehat{\mu}$.
\item For every $\varphi \in C_b(\RR)$, $\int \varphi(x)\mu_\sigma(dx) \longrightarrow \int \varphi(x)\mu(dx)$ as $\sigma \to 0$.
\end{enumerate}

In order to establish (1), we use the formula of the characteristic function for gaussian distribution to write, for every $x \in \RR$,
\[ \sigma \sqrt{2\pi}\, g_\sigma(x) = \exp\left(-\frac{x^2}{2\sigma^2}\right) = \int_{\RR} e^{i\xi x} g_{1/\sigma}(\xi)\, d\xi. \]

It follows that
\begin{align*}
f_\sigma(x) 
    &= \int_{\RR} g_\sigma(x - y)\,\mu(dy) \\
    &= (\sigma \sqrt{2\pi})^{-1} \int_{\RR} \left( \int_{\RR} e^{i\xi(x-y)} g_{1/\sigma}(\xi)\, d\xi \right)\mu(dy) \\
    &= (\sigma \sqrt{2\pi})^{-1} \int_{\RR} e^{i\xi x} g_{1/\sigma}(\xi)
    \left( \int_{\RR} e^{-i\xi y}\mu(dy) \right)d\xi \\
    &= (\sigma \sqrt{2\pi})^{-1} \int_{\RR} e^{i\xi x} g_{1/\sigma}(\xi)\,\widehat{\mu}(-\xi)\, d\xi.
\end{align*}

In the penultimate equality, we used the Fubini-Lebesgue theorem \ref{thm:Fubini_Lebesgue}. The application of this theorem is easily justified since $\mu$ is a probability measure and the function $g_{1/\sigma}$ is integrable with respect to Lebesgue measure. The formula of the last display shows that $f_\sigma$ (and hence $\mu_\sigma$) is determined by $\widehat{\mu}$.
As for (2), we start by writing, for every $\varphi \in C_b(\RR)$,
\begin{align*}
\int \varphi(x)\mu_\sigma(dx)
    &= \int \varphi(x)\left(\int g_\sigma(y - x)\mu(dy)\right)dx \\
    &= \int \left( \int g_\sigma(y - x)\varphi(x)\,dx \right)\mu(dy) \\
    &= \int g_\sigma * \varphi(y)\,\mu(dy).
\end{align*}
where we have again applied the Fubini-Lebesgue theorem, with the same justification as above, also the last equality use the convolution that we will see in \ref{sec:measure_product_and_independence}. Then, we note that the properties
\[\begin{cases}
\int g_\sigma(x)\,dx = 1, \\
\lim_{\sigma \to 0} \int_{\{|x| > \varepsilon\}} g_\sigma(x)\,dx = 0, \quad \forall \varepsilon > 0.
\end{cases}\]

imply that, for every $y \in \RR$,
\[ \lim_{\sigma \to 0} g_\sigma * \varphi(y) = \varphi(y) \]
thanks to Proposition \ref{prop:dirac_convolution} (1) that we will see in the section measure product and independence \ref{sec:measure_product_and_independence}. By dominated convergence, using the fact that $|g_\sigma * \varphi| \leq \sup\{|\varphi(x)| : x \in \RR\}$, we get 
\[ \lim_{\sigma \to 0} \int \varphi(x)\mu_\sigma(dx) = \lim_{\sigma \to 0} \int g_\sigma * \varphi(y)\mu(dy) = \int \varphi(x)\mu(dx), \]
which completes the proof of (2) and of the case $d = 1$ of the theorem. \\ 
The proof in the general case is similar. We now use the auxiliary functions
\[
g_\sigma^{(d)}(x_1,\ldots,x_d) = \prod_{j=1}^d g_\sigma(x_j)
\]
noting that, for $\xi = (\xi_1,\ldots,\xi_d) \in \RR^d$,
\[
\int_{\RR^d} g_\sigma^{(d)}(x) e^{i\xi \cdot x}\, dx
= \prod_{j=1}^d \int_{\RR} g_\sigma(x_j)e^{i\xi_j x_j}\, dx_j
= (2\pi/\sigma^2)^{d/2} g_{1/\sigma}^{(d)}(\xi).
\]

The arguments of the case $d = 1$ then apply with obvious changes.
}

\proppr{taylor_expansion_characteristic_function}{
Let $X = (X_1,\ldots,X_d)$ be a random variable with values in $\RR^d$. Assume that $\EE[(X_j)^2] < \infty$ for every $j \in \{1,\ldots,d\}$. Then, $\Phi_X$ is twice continuously differentiable and
\[
\Phi_X(\xi) = 1 + i \sum_{j=1}^d \xi_j \EE[X_j]
- \frac{1}{2} \sum_{j=1}^d \sum_{k=1}^d \xi_j \xi_k \EE[X_j X_k]
+ o(|\xi|^2)
\]
when $\xi = (\xi_1,\ldots,\xi_d)$ tends to $0$.
}{
By differentiating under the integral sign (easy application of dominance convergence theorem), we get, for every $1 \leq j \leq d$,
\[ \frac{\partial \Phi_X}{\partial \xi_j}(\xi) = i\, \EE[X_j e^{i\xi \cdot X}]. \]

The justification is easy since $|i X_j e^{i\xi \cdot X}| = |X_j|$ and $X_j \in L^2(\Omega,\A,\PP) \subset L^1(\Omega,\A,\PP)$. Similarly, for every $j,k \in \{1,\ldots,d\}$, the fact that $X_j X_k \in L^1(\Omega,\A,\PP)$ allows us to differentiate once more and to get
\[ \frac{\partial^2 \Phi_X}{\partial \xi_j \partial \xi_k}(\xi) = - \EE[X_j X_k e^{i\xi \cdot X}]. \]

Moreover, the dominated convergence theorem ensures that the quantity $\EE[X_j X_k e^{i\xi \cdot X}]$ is a continuous function of $\xi$. We have thus proved that $\Phi_X$ is twice continuously differentiable.

Finally, the last assertion of the proposition is just Taylor’s expansion of $\Phi_X$ at order two at the origin.
}

\subsection{Laplace transform}
For a random variable with values in $\RR_+$, one often uses the Laplace transform instead of the characteristic function.

\defnr{Laplace transform}{laplace_transform}{
Let $X$ be a nonnegative random variable. The Laplace transform of (the law of) $X$ is the function $L_X$ defined on $\RR_+$ by
\[ L_X(\lambda) := \EE[e^{-\lambda X}] = \int_{\RR_+} e^{-\lambda x} \,\PP_X(dx). \]
}

The dominated convergence theorem shows that $L_X$ is continuous on $[0,\infty)$. An application of the theorem of differentiation under the integral [which is just easy result of convergence theorem] sign shows that $L_X$ is infinitely differentiable on $(0,\infty)$. The function $L_X$ is also convex (as a consequence of the fact that the functions $\lambda \mapsto e^{-\lambda x}$ are convex, for every $x \geq 0$) and thus has a (possibly infinite) right derivative at $0$, which is given by
\[ L_X'(0) = -\lim_{\lambda \downarrow 0} \EE\left[\frac{1 - e^{-\lambda X}}{\lambda}\right] = -\EE[X] \]
where the second equality follows from the monotone convergence theorem.

\thmpr{Laplace transform determines the law}{law_determined_by_laplace_transform}{
If $X$ is a nonnegative random variable, the Laplace transform of $X$ determines the law $\PP_X$ of $X$.
}{
To prove this theorem, the reader may follow the steps below:
\begin{enumerate}
    \item Show that equality of the Laplace transforms implies equality of expectations for the exponential functions \(x \mapsto e^{-\lambda x}\), and hence for all their finite linear combinations.
    
    \item Apply the Stone--Weierstrass theorem to uniformly approximate every function in \(C([0,\infty])\) by elements of the linear span of the exponential functions.
    
    \item Pass to the limit using the uniform convergence obtained in the previous step to extend the equality of expectations to all continuous functions.
    
    \item Conclude that the induced probability measures coincide by the measure-determining property of continuous functions with compact support.
\end{enumerate}
}

Finally, in the case of integer-valued random variables, one often uses generating functions.

\defnr{Generating function}{Generating_function}{
Let $X$ be a random variable with values in $\NN$. The generating function of $X$ is the function $g_X : [0,1] \to [0,1]$ defined by
\[
g_X(r) := \EE[r^X] = \sum_{n=0}^{\infty} \PP(X = n)\, r^n.
\]
}

Note that $g_X(0) = \PP(X = 0)$ and $g_X(1) = 1$. The function $g_X$ is continuous on $[0,1]$. Since the radius of convergence of the series in the definition is at least $1$, the function $g_X$ is infinitely differentiable on $[0,1)$. Moreover, it is obvious that $g_X$ determines $\PP_X$, since the quantities $\PP(X = n)$ appear in the Taylor expansion of $g_X$ at the origin.

The derivative of $g_X$ is
\[ g_X'(r) = \sum_{n=1}^{\infty} n\, \PP(X = n)\, r^{n-1} \]
for $r \in [0,1)$. It follows that $g_X$ has a (possibly infinite) left-derivative at $r = 1$, and
\[ g_X'(1) = \sum_{n=1}^{\infty} n\, \PP(X = n) = \EE[X]. \]

More generally, for every integer $p \geq 1$, if $g_X^{(p)}$ denotes the $p$-th derivative of $g_X$,
\[\lim_{r \uparrow 1} g_X^{(p)}(r) = \EE[X(X-1)\cdots(X-p+1)]\]
which shows how to recover all moments of $X$ from the knowledge of its generating function.

\section{Lp Spaces} \label{sec:Lp spaces}
This section is mainly devoted to the study of the Lebesgue space $\L^p$ of random variables whose $p$-th power of the absolute value is having expectaiton (integrable) on a given probability measure. The fundamental inequalities of Hölder, Minkowski and Jensen provide essential ingredients for this study. We investigate the Banach space structure of $\L^p$ , and in the special case $p = 2$, the Hilbert space structure of $\L^2$, which has important applications in probability theory. 

Density theorems showing that any function of $\L^p$ can be well approximated by \textit{smoother} functions play an important role in many analytic developments. As an application of the Hilbert space structure of $L^2$ , we establish the Radon-Nikodym theorem, which allows one to decompose an arbitrary measure as the sum of a measure having a density with respect to the reference measure and a singular measure. The Radon-Nikodym theorem will be a crucial ingredient of the theory of conditional expectations.


\subsection{Definition of the space}
Throughout this section, we consider a probability space $(\probspace)$. For a real $p \geq 1$, we let $\L^p_{\RR^d}(\probspace)$ denote the space of all random variables $\bX: \Omega  \mapsto \RR^d$ such that:
\[ \int|\bX|^p \mathrm{~d} \PP<\infty  \]
We also introduce the space $\L^{\infty}_{\RR^d}(\probspace)$ of all random variables $\bX: \Omega \rightarrow \RR$ such that there exists a constant $C \in \RR_{+}$with
\[ |\bX| \leq C, \PP \text { a.e. } \]

For every $p\in[1,\infty]$, we define an equivalence relation on the set of
$p$-integrable random variables by setting : 
\[ \bX \equiv \bY \quad \text{if and only if} \quad \bX=\bY,\ \PP\text{-a.s.} \]

We then consider the quotient space
\[ L^p_{\RR^d}(\probspace)= \L^p_{\RR^d}(\probspace)/\equiv, \]

An element of $L^p_{\RR^d}(\probspace)$ is thus an equivalence class of random variables that are equal almost surely. In fact, $\conj{X}\in\L^p_{\RR^d}(\probspace)$ is defined as : 
\[ \conj{X} = \left\{ Y \in \L^p_{\RR^d}(\Omega,\A,\PP)\quad \text{such that } \bX = \bY \quad \PP-a.s. \right\}\] 
In fact we can define operation 


In what follows, we will almost always abuse notation and identify an element of $L^p_{\RR^d}(\probspace)$ with one of its representatives: when we speak of a random variable in $L^p_{\RR^d}(\probspace)$, we mean that a particular representative has been chosen, and all subsequent considerations do not depend on this choice.

We will often write $L^p(\probspace)$, $L^p(\PP)$, $L^p_{\RR^d}$ or simply $L^p$, instead of $L^p_{\RR^d}(\probspace)$ when there is no risk of confusion. Notice that the space $L^0$ is the space of all random variables, the space $L^1$ is precisely the space of integrable random variables, and that the space $L^2$ is the space of square-integrable random variables, which is of particular importance in probability theory.

For every random variable $\bX$ on $(\probspace)$ and every $p\in[1,\infty)$, we set : 
\[ \|\bX\|_p = \left(\EE\!\left[|\bX|^p\right]\right)^{1/p}, \]
with the convention $\infty^{1/p}=\infty$, and
\[ \|\bX\|_\infty = \inf\{C\in[0,\infty] : |\bX|\le C,\ \PP\text{-a.s.}\}. \]

These definitions are such that $|\bX|\le\|\bX\|_\infty$ $\PP$-a.s., and $\|\bX\|_\infty$ is the smallest number in $[0,\infty]$ with this property. We observe that, if $X$ and $Y$ are two random variables such that $X=Y$ $\PP$-a.s., then $\|\bX\|_p=\|\bY\|_p$. Hence $\|\bX\|_p$ is well defined for $\bX\in L^p(\probspace)$.

It is worth noting that when dealing with random variables taking values in $\RR^d$, the absolute value can be replaced by the Euclidean norm, indeed for a for a vector  $|\bx|$ the euclidian norm is $\sqrt{x_1^1+...+x_d^2}$.



\subsection{Principale inequalities}
\defnr{Conjugate exponents}{conjugate_exponents}{
    Let $p,q\in[1,\infty]$. We say that $p$ and $q$ are \emph{conjugate exponents} if
\[ \frac{1}{p}+\frac{1}{q}=1.\]
In particular, $p=1$ and $q=\infty$
}
\thmpr{Holder Inequality}{Holder_inequality}{
Let $p$ and $q$ be conjugate exponents. Then, if $\bX$ and $\bY$ are two random variables on $(\probspace)$,
\[ \EE\!\left[|\bX\bY|\right] \le \|\bX\|_p\,\|\bY\|_q. \]
In particular, $\bX\bY\in L^1(\probspace)$ whenever $\bX\in L^p(\probspace)$ and $\bY\in L^q(\probspace)$.
}{
If $\|\bX\|_p=0$, then $\bX=0$ $\PP$-a.s., which implies $\EE[|\bX\bY|]=0$, and the inequality is trivial. We may therefore assume that $\|\bX\|_p>0$ and $\|\bY\|_q>0$. Without loss of generality, we may also assume that $\bX\in L^p(\probspace)$ and $\bY\in L^q(\probspace)$ (otherwise $\|\bX\|_p\|\bY\|_q=\infty$ and there is nothing to prove).

The case $p=1$ and $q=\infty$ is immediate, since $|\bX\bY|\le \|\bY\|_\infty |\bX|$ $\PP$-a.s., which implies :
\[ \EE[|\bX\bY|] \le \|\bY\|_\infty\,\EE[|\bX|]  = \|Y\|_\infty\,\|\bX\|_1. \]

In what follows, we assume that $1<p<\infty$ (and thus $1<q<\infty$). Let $\alpha\in(0,1)$. Then, for every $x\in\RR_+$,
\[ x^\alpha-\alpha x \le 1-\alpha. \]
Indeed, define $\varphi_\alpha(x)=x^\alpha-\alpha x$ for $x\ge0$. For $x>0$,
\[ \varphi_\alpha'(x)=\alpha(x^{\alpha-1}-1), \]
so that $\varphi_\alpha'(x)>0$ for $x\in(0,1)$ and $\varphi_\alpha'(x)<0$ for $x\in(1,\infty)$. Hence $\varphi_\alpha$ attains its maximum at $x=1$, which yields the inequality above.

Applying this inequality to $x=u/v$, where $u\ge0$ and $v>0$, we obtain
\[ u^\alpha v^{1-\alpha} \le \alpha u + (1-\alpha)v, \]
and the inequality still holds when $v=0$. We now take $\alpha=1/p$ (so that $1-\alpha=1/q$) and set :
\[ u=\frac{|\bX(\omega)|^p}{\|\bX\|_p^p}, \qquad v=\frac{|\bY(\omega)|^q}{\|\bY\|_q^q}.\]
This yields
\[ \frac{|\bX(\omega)Y(\omega)|}{\|\bX\|_p\|\bY\|_q} \le \frac{1}{p}\frac{|\bX(\omega)|^p}{\|\bX\|_p^p} + \frac{1}{q}\frac{|Y(\omega)|^q}{\|\bY\|_q^q}, \quad \PP\text{-a.s.}
\]

Taking expectations on both sides gives
\[ \frac{1}{\|\bX\|_p\|\bY\|_q}\,\EE[|\bX\bY|] \le \frac{1}{p}+\frac{1}{q} = 1, \]
which completes the proof.
}

\thmpr{Cauchy-Schwartz Inequality}{CauchySchwartz_inequality}{
If $\bX$ and $\bY$ are two random variables on $(\probspace)$, then
\[ \EE\!\left[|\bX\bY|\right] \le \left(\EE\!\left[|\bX|^2\right]\right)^{1/2} \left(\EE\!\left[|\bY|^2\right]\right)^{1/2} = \|\bX\|_2\,\|\bY\|_2. \]
}{
    Just a special case of Holder Inequality with $p=2$
}

\proppr{Lp_spaces_inclusion}{
Suppose that $\PP$ is a probability measure and let $p>1$. Then, for any random variable $\bX\in L^p(\probspace)$,
\[ \|\bX\|_1 \le \|\bX\|_p. \]
Consequently, $L^p(\probspace)\subset L^1(\probspace)$ for every $p\in(1,\infty]$.
More generally, for any $1\le r<r'\le\infty$ and any random variable
$\bX\in L^{r'}(\probspace)$,
\[ \|\bX\|_r \le \|\bX\|_{r'}, \]
and thus $L^{r'}(\probspace)\subset L^r(\probspace)$.
}{
The inequality $\|\bX\|_1 \le \|\bX\|_p$ follows directly from Hölder's inequality by taking $\bY=1$. For the second bound in the corollary, it suffices to replace $\bX$ by $|\bX|^r$ and to take $p = r'/r$.
}   


\thmpr{Jensen Inequality}{Jensen_inequality}{
Suppose that $\PP$ is a probability measure on $(\measspace)$, and let $\varphi:\RR\to\RR_+$ be a convex function. Then, for every integrable real-valued random variable $X\in L^1_{\RR}(\probspace)$,
\[ \EE\!\left[\varphi(X)\right] \ge \varphi\!\left(\EE[X]\right). \]
\textbf{remark :} The expectation $\EE[\varphi(X)]$ is well defined, since $\varphi(X)$ is a nonnegative measurable random variable.
}{
Set
\[ \E_\varphi = \{(a,b)\in\RR^2 : \forall x\in\RR,\ \varphi(x)\ge ax+b\}. \]
Elementary properties of convex functions show that
\[ \forall x\in\RR,\quad \varphi(x)=\sup_{(a,b)\in\E_\varphi}(ax+b). \]

Since $\varphi(X)\ge aX+b$ almost surely for every $(a,b)\in\E_\varphi$, we obtain
\[ \EE[\varphi(X)] \ge \sup_{(a,b)\in\E_\varphi} \EE[aX+b] = \sup_{(a,b)\in\E_\varphi} \bigl(a\,\EE[X]+b\bigr) = \varphi\!\left(\EE[X]\right). \]   
}


\thmpr{Minkowski's Inequality}{Minkowski_inequality}{
Let $p\in[1,\infty]$, and let $\bX,\bY \in L^p(\probspace)$. Then :
\[ \bX+\bY\in L^p(\probspace) \quad \text{and} \quad \|\bX+\bY\|_p \le \|\bX\|_p + \|\bY\|_p. \]
}{
The cases $p=1$ and $p=\infty$ are immediate, using the inequality $|\bX+\bY|\le |\bX|+|\bY|$. We therefore assume that $1<p<\infty$.

Writing
\[ |\bX+\bY|^p \le (|\bX|+|\bY|)^p \le \bigl(2\max(|\bX|,|\bY|)\bigr)^p \le 2^p\bigl(|\bX|^p+|\bY|^p\bigr), \]
we see that $\EE[|\bX+\bY|^p]<\infty$, and hence $\bX+\bY\in L^p$.

Next, observe that
\[ |\bX+\bY|^p = |\bX+\bY|\cdot |\bX+\bY|^{p-1} \le |\bX|\,|\bX+\bY|^{p-1} + |\bY|\,|\bX+\bY|^{p-1}. \]
Taking expectations yields
\[ \EE[|\bX+\bY|^p] \le \EE\!\left[|\bX|\,|\bX+\bY|^{p-1}\right] + \EE\!\left[|\bY|\,|\bX+\bY|^{p-1}\right]. \]

Applying Hölder's inequality with conjugate exponents $p$ and $q=p/(p-1)$, we obtain
\[ \EE[|\bX+\bY|^p] \le \|\bX\|_p \left(\EE[|\bX+\bY|^p]\right)^{(p-1)/p} + \|\bY\|_p \left(\EE[|\bX+\bY|^p]\right)^{(p-1)/p}. \]

If $\EE[|\bX+\bY|^p]=0$, the inequality is trivial. Otherwise, dividing both sides by $\left(\EE[|\bX+\bY|^p]\right)^{(p-1)/p}$ yields
\[ \|\bX+\bY\|_p \le \|\bX\|_p + \|\bY\|_p, \]
which completes the proof.
}

\subsection{Topological structure of Lp-space}



\proppr{L0_complete}{
The formula
\[ d(X,Y) = \EE\bigl[\,|X - Y| \wedge 1\,\bigr] \]
defines a distance on $L^0_{\RR^d}(\Omega,\A,\PP)$, and the space $L^0_{\RR^d}(\Omega,\A,\PP)$ is complete for the distance $d$.
}{
Let's prove that the space $L^0_{\RR^d}(\Omega,\A,\PP)$ is complete for the distance $d$. So let $(X_n)_{n \in \NN}$ be a Cauchy sequence for the distance $d$. We can then find a subsequence $Y_k = X_{n_k}$, $k \in \NN$, such that for every $k \ge 1$,
\[d(Y_k,Y_{k+1}) \le 2^{-k}.\]

Then,
\[ \EE\Biggl[\sum_{k=1}^{\infty} \bigl(|Y_{k+1}-Y_k| \wedge 1\bigr)\Biggr] = \sum_{k=1}^{\infty} d(Y_k,Y_{k+1}) < \infty, \]
which implies that $\sum_{k=1}^{\infty} (|Y_{k+1}-Y_k| \wedge 1) < \infty$ almost surely, and therefore also $\sum_{k=1}^{\infty} |Y_{k+1}-Y_k| < \infty$ almost surely (there are almost surely only finitely many values of $k$ such that $|Y_{k+1}-Y_k| \ge 1$). We then define a random variable in $L^0_{\RR^d}(\Omega,\A,\PP)$ by setting
\[X = Y_1 + \sum_{k=1}^{\infty} (Y_{k+1}-Y_k),\]
on the event where the series converges absolutely, and $X=0$ on the complementary event, which has zero probability. By construction, the sequence $(Y_k)_{k \in \NN}$ converges almost surely to $X$, and this implies
\[ d(Y_k,X) = \EE\bigl[\,|Y_k - X| \wedge 1\,\bigr] \longrightarrow 0 \quad \text{as } k \to \infty, \]
by an application of the dominated convergence theorem. Hence the sequence $(Y_k)_{k \in \NN}$ converges to $X$ for the distance $d$, and since the sequence $(X_n)_{n \in \NN}$ is Cauchy and has a convergent subsequence, it must also converge.
}


\thmpr{Riesz Theorem}{Riesz}{ 
For every $p\in[1,\infty]$, the space $L^p(\probspace)$ equipped with the norm $\bX\mapsto \|\bX\|_p$ is a Banach space, that is, a complete normed linear space.
}{
    Can be found in  \parencite[p.~68]{Probability}
}

\thmpr{The Hilbert Space L2}{L2_Hilbert_space}{
The space $L^2_\RR(\probspace)$ equipped with the scalar product
\[ \langle X, Y \rangle = \int XY \, d\PP = \EE(XY) \]
is a real Hilbert space.
}{ The Cauchy-Schwarz inequality shows that, if $X,Y\in L^2(\probspace)$, then $XY$ is integrable and hence $\langle X,Y\rangle$ is well defined. It is then clear that the mapping $(X,Y)\mapsto \langle X,Y\rangle$ defines an inner product on $L^2(\probspace)$, and that $\langle X,X\rangle=\|X\|_2^2$. Finally, the fact that $\bigl(L^2(\probspace),\|\cdot\|_2\bigr)$ is complete is a special case of the Riesz theorem \ref{thm:Riesz}.
}
Classical results from the theory of Hilbert spaces [see appendix \ref{apen:Functional_Analysis}] can therefore be applied to $L^2(\probspace)$. In particular, if $\Phi : L^2_{\RR}(\probspace)\to\RR$ is a continuous linear functional, then there exists a unique element $Y\in L^2_{\RR^d}(\probspace)$ such that
\[ \Phi(X)=\langle X,Y\rangle \quad\text{for every } X\in L^2_{\RR^d}(\probspace), \]
where $\langle X,Y\rangle=\EE[XY]$; .

\subsection{The Radon-Nikodym Theorem}
\defnr{Measures continuity}{Measures_continuity}{
Let $\PP$ and $\QQ$ be two probability measures on $(\measspace)$. We say that:
\begin{enumerate}
\item $\QQ$ is \emph{absolutely continuous} with respect to
$\PP$ (notation $\QQ \ll \PP$) if
\[ \forall A\in\A,\quad \PP(A)=0 \;\Rightarrow\; \QQ(A)=0. \]

\item $\QQ$ is \emph{singular} with respect to $\PP$ (notation $\QQ \perp \PP$) if there exists $\Omega_0\in\A$ such that
\[ \PP(\Omega_0)=0 \quad\text{and}\quad \QQ(\Omega_0^c)=0. \]

\item $\PP$ and $\QQ$ are said to be \emph{equivalent} (notation $\PP \sim \QQ$) if
\[ \PP \ll \QQ \quad\text{and}\quad \QQ \ll \PP. \] 
\end{enumerate}
}
Now we present the theorem of Radon-Nikodym : 
\thmpr{Radon-Nikodym Theorem}{RadonNikodym}{
Suppose that $\PP$ is a $\sigma$-finite measure\footnote{A measure $\PP$ on $(\Omega,\A)$ is called \emph{finite} if $\PP(\Omega)<\infty$. It is called \emph{$\sigma$-finite} if there exists a sequence of measurable sets $(A_n)_{n\ge1}$ such that $\Omega=\bigcup_{n=1}^{\infty} A_n$ and $\PP(A_n)<\infty$ for all $n$.} on $(\measspace)$, and let $\QQ$ be another $\sigma$-finite measure on $(\measspace)$. Then there exists a unique pair $(\QQ_a,\QQ_s)$ of $\sigma$-finite measures on $(\measspace)$ such that:
\begin{enumerate}
\item $\QQ = \QQ_a + \QQ_s$,
\item $\QQ_a \ll \PP$ and $\QQ_s \perp \PP$.
\end{enumerate}
Moreover, there exists a unique\footnote{Uniqueness is in the sense that if $\tilde g$ is another measurable function satisfying the same properties, then $\tilde g = g$ $\PP$-a.s.} nonnegative measurable function $g:\Omega\to\RR_+$ such that
\[ d\QQ_a = g\cdot d\PP, \]
meaning that
\[ \forall A\in\A,\qquad \QQ_a(A)=\int_A g\,d\PP. \]
}{
    Could be found in \parencite[p.~76]{Probability} proven in the cadre of measure theory.
}
\textbf{Remark:} A useful formula one can use in \textit{change of measures} in Radon-Nikodym calculations is this : 
\[ d\QQ = g\cdot d\PP \quad\quad \textit{then} \quad\quad \int_A f d\QQ =  \int_A f\cdot g\cdot d\PP \]

\thmr{Radon-Nikodym Theorem in Probability}{RadonNikodym_prob}{
    Suppose that $\PP$ and $\QQ$ are two probability measures such that $\QQ \ll \PP$ then :\\
    There exists a unique\footnote{Uniqueness is in the sense that if $\tilde X$ is another random variable satisfying the same properties, then $\tilde X = X$ $\PP$-a.s.} nonnegative random variable $X:\Omega\to\RR_+$ such that
    \[ d\QQ = X\cdot d\PP, \]
    Furthermore, if $\PP \ll \QQ$ then $1/X$ satisfies :
    \[ d\PP = \frac{1}{X}\cdot d\QQ, \]
    
}


\proppr{equivalent_definition_absolute_continuity}{
Let $\PP$ be a $\sigma$-finite measure on $(\measspace)$, and let $\QQ$ be a finite measure on $(\measspace)$. The following conditions are equivalent:
\begin{enumerate}
\item $\QQ \ll \PP$;
\item For every $\varepsilon>0$, there exists $\delta>0$ such that, for every $A\in\A$,
\[ \PP(A)<\delta \;\Longrightarrow\; \QQ(A)<\varepsilon. \]
\end{enumerate}
}{
The implication $(\mathrm{ii}) \Rightarrow (\mathrm{i})$ is immediate.
Conversely, suppose that $\QQ \ll \PP$. By the Radon-Nikodym theorem \ref{thm:RadonNikodym}, there exists a nonnegative measurable function $g:\Omega\to\RR_+$ such that
\[ d\QQ = g\cdot d\PP. \]
Note that
\[ \int_\Omega g\,d\PP = \QQ(\Omega) < \infty. \]
The dominated convergence theorem then yields
\[ \int_\Omega \1_{\{g>n\}}\, g\, d\PP \conv{n} 0. \]
Given $\varepsilon>0$, we may therefore choose $M\in\NN$ large enough such that
\[ \int_\Omega \1_{\{g>M\}}\, g\, d\PP < \frac{\varepsilon}{2}. \]
Setting $\delta=\varepsilon/(2M)$, we obtain, for every $A\in\A$ such that $\PP(A)<\delta$,
\[ \QQ(A) = \int_A g\,d\PP \le \int_\Omega \1_{\{g>M\}}\, g\, d\PP + \int_{A\cap\{g\le M\}} g\,d\PP < \frac{\varepsilon}{2} + M\,\PP(A) \le \varepsilon. \]
This proves $(\mathrm{i}) \Rightarrow (\mathrm{ii})$ and completes the proof.
}


\end{document}